\chapter{Giardiasis (Giardia)}

\section*{Definition and Overview}
Giardiasis (giardia) is an intestinal infection caused by the microscopic parasite \textit{Giardia}, one of the most common causes of waterborne diarrhoeal disease in the world. The parasite lives in the small intestine and is passed in the faeces of infected people and animals. People become infected by swallowing giardia cysts, usually through contaminated drinking water, food, or hands. The infection can cause diarrhoea, stomach cramps, bloating, and nausea, but many people infected with giardia have no symptoms at all. Giardiasis is easily treated with prescription antiparasitic medicines, and good hygiene and safe water practices prevent it. This chapter focuses on the practical aspects of recognising, treating, and preventing giardia infection.

\section*{Epidemiology}
Giardia is found worldwide and infects an estimated 200 million people each year. It is a common cause of diarrhoea in children, a well-known hazard for campers and hikers drinking untreated surface water, and an important infection in travellers returning from areas with poor sanitation. In Australia, giardiasis is a notifiable disease, with thousands of cases reported each year, and outbreaks are frequently traced to contaminated water supplies, swimming pools, and childcare centres. Young children have the highest rates of infection, and households and carers of infected children are at increased risk.

\section*{Aetiology and Pathophysiology}
Giardia infection follows a well-understood cycle involving the hardy cyst form and the disease-causing trophozoite form of the parasite.

\begin{itemize}
    \item \textbf{The cyst:} giardia cysts are passed in stool and can survive for months in cold water and withstand routine chlorination
    \item \textbf{Ingestion:} infection occurs when cysts are swallowed in contaminated water, food, or via hands that have touched contaminated surfaces
    \item \textbf{Excystation:} in the small intestine, cysts break open to release trophozoites
    \item \textbf{Attachment:} trophozoites attach to the intestinal wall, damaging the microvilli and impairing digestion and absorption
    \item \textbf{Reproduction:} trophozoites multiply and eventually encyst again, producing new cysts that are shed in the stool
    \item \textbf{Spread:} poor hand hygiene, contaminated water sources, and faecal contamination of food perpetuate transmission
\end{itemize}

\section*{Clinical Features}
The incubation period is usually one to three weeks, and symptoms vary from none to severe.

\begin{itemize}
    \item Acute or chronic diarrhoea, often watery, greasy, and particularly foul-smelling
    \item Abdominal cramps, bloating, and excessive flatulence
    \item Nausea, poor appetite, and weight loss
    \item Fatigue and malaise
    \item Stools that float due to fat content
    \item Dehydration in severe or prolonged cases
    \item Many people, especially children, carry the parasite without any symptoms
    \item Chronic infection can cause intermittent symptoms for months
\end{itemize}

\section*{Differential Diagnosis}
Because giardia symptoms overlap with many other gut infections and conditions, laboratory confirmation is required.

\begin{itemize}
    \item Cryptosporidiosis, another waterborne parasite
    \item Bacterial gastroenteritis from campylobacter, salmonella, shigella, or \textit{E. coli}
    \item Viral gastroenteritis from norovirus, rotavirus, or adenovirus
    \item Amoebiasis and other intestinal parasites
    \item Coeliac disease and other malabsorptive disorders
    \item Irritable bowel syndrome
    \item Inflammatory bowel disease
    \item Food poisoning from toxins or contaminated food
\end{itemize}

\section*{When to Seek Medical Care}
Seek medical advice if diarrhoea lasts more than a few days, is severe, or is associated with dehydration, weight loss, or blood in the stool. Tell your doctor about recent travel, camping, or contact with anyone who has had giardia.

\textbf{Call 000 (or local emergency services) for:}
\begin{itemize}
    \item Signs of severe dehydration, including little or no urine, extreme thirst, lethargy, or confusion
    \item Bloody diarrhoea or severe abdominal pain
    \item Diarrhoea with persistent vomiting in a young child or baby
    \item Any person with diarrhoea who becomes weak, dizzy, or unable to drink
\end{itemize}

\section*{Diagnosis and Investigations}
Giardiasis is confirmed by laboratory testing of stool, because symptoms alone cannot distinguish it from other causes of gastroenteritis.

\begin{itemize}
    \item \textbf{Stool microscopy:} multiple samples over several days improve detection, as cysts are shed intermittently
    \item \textbf{Antigen testing:} sensitive rapid tests detect giardia proteins in stool
    \item \textbf{PCR testing:} detects giardia DNA and is highly sensitive
    \item \textbf{History:} exposure to untreated water, childcare, travel, and contacts helps guide and interpret testing
    \item \textbf{Response to treatment:} improvement with antiparasitic therapy supports the diagnosis
\end{itemize}

\section*{Management}
Giardiasis responds well to specific treatment, and supportive care manages symptoms while the medicine works.

\begin{itemize}
    \item \textbf{Antiparasitic medicines:} tinidazole or metronidazole are first-line treatments; nitazoxanide is an option for children
    \item \textbf{Fluids:} oral rehydration solution prevents and corrects dehydration
    \item \textbf{Dietary measures:} eat bland foods and temporarily reduce dairy and fatty foods if they worsen symptoms
    \item \textbf{Treatment of household contacts:} people with symptoms should be tested and treated to prevent re-infection
    \item \textbf{Repeat testing:} may be arranged after treatment in some settings, though it is not needed for everyone
    \item \textbf{Specialist care:} people with weakened immunity may need longer courses and follow-up
\end{itemize}

\section*{Complications}
\begin{itemize}
    \item Chronic diarrhoea and malabsorption with weight loss and vitamin deficiencies
    \item Dehydration, particularly in children and older people
    \item Lactose intolerance that may persist after the infection
    \item Failure to thrive in young children with chronic infection
    \item Post-infectious irritable bowel syndrome
    \item Prolonged or recurrent infection in people with immune deficiency
    \item Ongoing transmission to household and community contacts if hygiene is not addressed
\end{itemize}

\section*{Prognosis and Outcomes}
The vast majority of people with giardiasis recover fully. With appropriate antiparasitic treatment, symptoms usually resolve within one to two weeks, and untreated infections often clear within several weeks as the immune system responds. Some people experience lingering digestive symptoms, and a small proportion develop persistent irritable bowel-type symptoms. In healthy people, giardiasis rarely causes lasting harm, and repeat infections are preventable with good hygiene and safe water practices.

\section*{Prevention and Self-Care}
\begin{itemize}
    \item Wash hands thoroughly with soap and water after using the toilet, changing nappies, and before eating or preparing food
    \item Treat water from rivers, lakes, and streams by boiling, filtering, or using certified purification tablets
    \item Avoid swallowing pool, lake, or river water when swimming
    \item Wash raw fruit and vegetables before eating
    \item Keep children with diarrhoea home from childcare until symptoms have settled
    \item Disinfect surfaces and toilet areas used by an infected person
    \item Avoid preparing food for others while you have diarrhoea
    \item Drink plenty of fluids if infected, and follow the treatment prescribed by your doctor
\end{itemize}

\section*{Sources and Further Reading}
The following reputable sources provide further information for healthcare students and patients seeking reliable, current advice about giardia infection.

\begin{itemize}
    \item Healthdirect Australia. \textit{Giardia (giardiasis).} https://www.healthdirect.gov.au/
    \item Australian Government Department of Health and Aged Care. \textit{Giardia fact sheet.} https://www.health.gov.au/
    \item Queensland Health. \textit{Giardiasis fact sheet.} https://www.health.qld.gov.au/
    \item Centers for Disease Control and Prevention. \textit{Giardia: general information.} https://www.cdc.gov/
    \item World Health Organization. \textit{Giardiasis fact sheet.} https://www.who.int/
    \item Mayo Clinic. \textit{Giardia infection (giardiasis).} https://www.mayoclinic.org/
\end{itemize}
